\documentclass[oneside]{article}
\usepackage[backend=biber,natbib=true,style=authoryear]{biblatex}
\addbibresource{/home/nguyen/1_NQBH/reference/bib.bib}
\usepackage[utf8]{inputenc}
\usepackage{float}
\usepackage{graphicx}
\usepackage[colorlinks=true,linkcolor=blue,urlcolor=red,citecolor=magenta]{hyperref}
\usepackage{amsmath,amssymb,amsthm,mathtools}
\allowdisplaybreaks
\numberwithin{equation}{section}
\newtheorem{assumption}{Assumption}[section]
\newtheorem{lemma}{Lemma}[section]
\newtheorem{corollary}{Corollary}[section]
\newtheorem{definition}{Definition}[section]
\newtheorem{proposition}{Proposition}[section]
\newtheorem{theorem}{Theorem}[section]
\newtheorem{notation}{Notation}[section]
\newtheorem{remark}{Remark}[section]
\newtheorem{example}{Example}[section]
\newtheorem{ques}{Question}[section]
\newtheorem{problem}{Problem}[section]
\newtheorem{conjecture}{Conjecture}[section]
\usepackage[left=0.5in,right=0.5in,top=1.5cm,bottom=1.5cm]{geometry}
\usepackage{fancyhdr}
\pagestyle{fancy}
\fancyhf{}
\lhead{\small \textsc{Sect.} ~\thesection}
\rhead{\small \nouppercase{\leftmark}} % \nouppercase !
\renewcommand{\sectionmark}[1]{\markboth{#1}{}}
\cfoot{\thepage}
\def\labelitemii{$\circ$}

\title{Collins. Unlocking German with Paul Noble: Your Key to Language Success. 2018}
\author{Collector: Nguyen Quan Ba Hong}
\date{\today}

\begin{document}
\maketitle
\setcounter{secnumdepth}{6}
\setcounter{tocdepth}{6}
\tableofcontents

%------------------------------------------------------------------------------%

\section*{Introduction: Did you know you already speak German?}
Did you know you already speak German?

That you speak it every day?

That you read and write it every day?

That you use it with your friends, with your family, at work, down the post office - even in the shower when you read the label on the shampoo bottle?

%
\textbf{Were you aware of that fact?}

%
\textbf{Well, even if you weren't, it's nevertheless true.}

%
Of course, you might not have realized at the time that what you were reading/saying/writing was actually German but Paul can prove to you that it was.

Just take a look at these words below:
\begin{example}
    winter, name, warm, baby, land, blind, wild, arm, hand, finger, blond, hammer, horn, ring, wind, butter, bitter, sand, wolf, hunger, rose, nest, gold
\end{example}

\section{Can you come over tonight?}

\section{I wouldn't like to do it now because I'm very busy.}

\section{I bought the ticket but I didn't see the film.}

\section{I gave the money to the taxi driver.}

\section{My father's car is old \& shabby. (part 1)}

\section{The car my's father's is old \& shabby. Erm, what? (part 2)}

\section{We're lucky that the weather's so good. (part 1)}

\section{We're lucky that the weather's so good. (part 2)}

\section{Do you think I'm oblivious to what's going on around me? (part 1)}

\section{Do you think I'm oblivious to what's going on around me? (part 2)}

%------------------------------------------------------------------------------%

\printbibliography[heading=bibintoc]
\end{document}